% =====================================================================
% WitnessRAG explicado: guia de leitura (não é texto de artigo)
% Compilar: latexmk -pdf witnessrag-teoria.tex
% Descreve o controlador de prova como está no código em 25/09/2026
% (wrag/methods/witnessrag.py::_retrieve_proof e wrag/witness/*).
% =====================================================================
\documentclass[11pt]{article}
\usepackage[a4paper,margin=2.4cm]{geometry}
\usepackage[T1]{fontenc}
\usepackage[utf8]{inputenc}
\usepackage{lmodern}
\usepackage[brazilian]{babel}
\usepackage{amsmath,amssymb,mathtools}
\usepackage{microtype}
\usepackage{booktabs,tabularx,array,multirow}
\usepackage{enumitem}
\usepackage{xcolor}
\usepackage{tikz}
\usetikzlibrary{arrows.meta,positioning,calc,fit,backgrounds,shapes.geometric,decorations.pathreplacing}
\usepackage{pgfplots}
\pgfplotsset{compat=1.17}
\usepackage[most]{tcolorbox}
\usepackage[hidelinks]{hyperref}
\usepackage{caption}
\usepackage{float}
\usepackage{needspace}
\usepackage{titlesec}
\captionsetup{font=small,labelfont=bf,skip=4pt}
\titleformat{\section}{\normalfont\Large\bfseries\color{blue!45!black}}{\thesection}{0.6em}{}
\titleformat{\subsection}{\normalfont\large\bfseries}{\thesubsection}{0.6em}{}
\setlength{\parskip}{4pt}
\setlength{\parindent}{0pt}
\setlist{leftmargin=*,itemsep=2pt,topsep=3pt}
\newcolumntype{Y}{>{\raggedright\arraybackslash}X}

\newcommand{\rel}[1]{\textsf{#1}}
\newcommand{\V}[1]{\mathord{?}\mathit{#1}}
\newcommand{\op}[1]{\textsc{#1}}
\newcommand{\sem}{\operatorname{sim}}
\newcommand{\prox}{\operatorname{prox}}
\DeclareMathOperator{\cost}{custo}
\DeclareMathOperator{\supp}{sup}

% Caixas
\newtcolorbox{formula}[1][]{breakable,enhanced,colback=blue!3,colframe=blue!45!black,
  fonttitle=\bfseries\small,title={Em fórmula#1},boxrule=0.6pt,arc=2pt,left=6pt,right=6pt,top=4pt,bottom=4pt}
\newtcolorbox{exemplo}[1][]{breakable,enhanced,colback=orange!4,colframe=orange!60!black,
  fonttitle=\bfseries\small,title={Exemplo#1},boxrule=0.6pt,arc=2pt,left=6pt,right=6pt,top=4pt,bottom=4pt}
\newtcolorbox{porque}[1][]{breakable,enhanced,colback=gray!6,colframe=gray!55,coltitle=black,
  fonttitle=\bfseries\small,title={Por que assim?#1},boxrule=0.6pt,arc=2pt,left=6pt,right=6pt,top=4pt,bottom=4pt}
\newtcolorbox{resumo}{enhanced,colback=green!4,colframe=green!40!black,
  fonttitle=\bfseries\small,title={Em uma frase},boxrule=0.6pt,arc=2pt,left=6pt,right=6pt,top=3pt,bottom=3pt}

\title{\textbf{Como o WitnessRAG funciona}\\[4pt]
\large Um guia passo a passo: a memória, o plano, a busca e a prova}
\author{}
\date{Estado do código em 25/09/2026}
\begin{document}
\maketitle
\thispagestyle{empty}

\section*{Como ler este guia}

Este texto explica o WitnessRAG para quem nunca viu o código: o que ele guarda,
como procura uma resposta e por que cada decisão foi tomada. Ele não é um
artigo: repete ideias quando isso ajuda e usa sempre o mesmo exemplo, uma
conversa curta entre Melanie e Caroline (no formato do benchmark LoCoMo).

Três tipos de caixa aparecem ao longo do texto:
\begin{itemize}
\item \textbf{Exemplo} (laranja): a ideia aplicada ao caso concreto, com números.
\item \textbf{Em fórmula} (azul): a mesma ideia escrita em matemática. Dá para
  pular na primeira leitura; o texto ao redor explica tudo sem ela.
\item \textbf{Por que assim?} (cinza): o motivo de uma escolha de desenho.
\end{itemize}
No fim há um glossário (Seção~\ref{sec:glossario}) e um mapa de onde cada
peça está no código (Seção~\ref{sec:codigo}).

\tableofcontents

% =====================================================================
\section{O problema}
\label{sec:problema}

Um assistente que conversa com alguém durante meses acumula centenas de
sessões. Quando chega uma pergunta (``Onde Melanie já acampou?''), ele não
consegue reler tudo. O que se faz na prática é dividir a conversa em
\emph{trechos}, escolher uns poucos (cinco, nos nossos experimentos) e
entregá-los a um modelo de linguagem, o \emph{leitor}, que escreve a resposta.

Tudo depende de escolher bem esses cinco trechos. O jeito mais comum é pegar
os trechos mais \emph{parecidos} com a pergunta, por palavras (BM25) e por
significado (embeddings). Isso funciona quando a resposta está em um trecho só
e usa palavras parecidas com as da pergunta. Falha em três situações comuns:

\begin{enumerate}
\item \textbf{A resposta está espalhada.} Melanie acampou nas montanhas, na
  praia e na floresta, e contou cada viagem em uma sessão diferente. Os cinco
  trechos mais parecidos com a pergunta podem trazer só duas das três.
\item \textbf{A resposta exige encadear fatos.} ``Para que cidade se mudou a
  mentora de Caroline?'' Um trecho diz que Beth é mentora de Caroline; outro,
  meses depois, diz que Beth se mudou para Boston. Esse segundo trecho
  \emph{nem cita Caroline}: ele só fica relevante depois que se sabe quem é a
  mentora. A similaridade com a pergunta não tem como achá-lo.
\item \textbf{A resposta depende de datas.} ``Quando Melanie acampou em
  junho?'' Os relatos de julho falam de acampamento tanto quanto o de junho.
  Para a similaridade, eles são igualmente bons.
\end{enumerate}

Nos três casos falta a mesma coisa: a busca por similaridade mede o quanto um
trecho se parece com a pergunta, mas não sabe dizer se os trechos escolhidos
\emph{bastam} para responder.

% =====================================================================
\section{A ideia: procurar uma prova}
\label{sec:ideia}

O WitnessRAG muda a pergunta que a busca faz. Em vez de ``quais trechos se
parecem com a pergunta?'', ele pergunta ``\emph{quais fatos provam a
resposta?}''.

Para isso, além do texto, a memória guarda \emph{fatos} curtos (``Beth é
mentora de Caroline'', ``Beth se mudou para Boston''), cada um com a data e a
fala de onde veio. Uma \textbf{testemunha} de uma resposta é o menor conjunto
de fatos do qual a resposta decorre. O nome vem da teoria de proveniência em
bancos de dados, onde a testemunha de um resultado é o conjunto de registros
que o produz.

\begin{exemplo}[: testemunhas]
\begin{itemize}
\item ``Para que cidade se mudou a mentora de Caroline?'' Resposta: Boston.
  Testemunha: \{Beth é mentora de Caroline; Beth se mudou para Boston\}. Dois
  fatos, e nenhum dos dois sozinho basta.
\item ``Onde Melanie já acampou?'' Três respostas, cada uma com sua
  testemunha de um fato só: \{Melanie acampou nas montanhas\},
  \{\dots na praia\}, \{\dots na floresta\}.
\end{itemize}
\end{exemplo}

A testemunha resolve os três problemas da seção anterior, porque dá duas
coisas que a similaridade não dá:
\begin{itemize}
\item \textbf{Um critério de parada.} Enquanto não há testemunha, a busca não
  terminou. Quando há, terminou.
\item \textbf{Uma lista do que o contexto precisa conter.} Cada fato da
  testemunha aponta para o trecho de onde veio. Esses trechos precisam estar
  entre os cinco que o leitor recebe. No exemplo da mentora, o trecho sobre
  Boston entra no contexto \emph{porque faz parte da prova}, mesmo sem se
  parecer com a pergunta.
\end{itemize}

\begin{resumo}
O WitnessRAG guarda a conversa também como fatos datados ligados entre si, e
para cada pergunta procura os poucos fatos que provam a resposta; os trechos
de onde vêm esses fatos entram no contexto do leitor.
\end{resumo}

% =====================================================================
\section{Visão geral}
\label{sec:visao}

O método tem duas fases (Figura~\ref{fig:visao}).

\begin{figure}[H]
\centering
\begin{tikzpicture}[
  box/.style={draw,rounded corners=3pt,thick,minimum height=10mm,minimum width=22mm,font=\sffamily\small,align=center},
  llm/.style={box,fill=orange!15}, det/.style={box,fill=blue!8},
  mem/.style={draw,rounded corners=3pt,thick,fill=gray!8,minimum height=11mm,align=center,font=\sffamily\small},
  arr/.style={-{Stealth[length=2.2mm]},thick}, lab/.style={font=\scriptsize,align=center}]
\node[font=\bfseries\small,anchor=west] at (-1.2,2.9) {Fase 1: uma vez por conversa};
\node[font=\small\itshape] (conv) at (-0.4,1.8) {conversa};
\node[llm,right=7mm of conv] (con) {Construir\\a memória};
\node[mem,right=7mm of con,minimum width=62mm] (m) {Memória: texto datado \;+\; grafo de fatos datados};
\draw[arr] (conv)--(con); \draw[arr] (con)--(m);
\node[font=\bfseries\small,anchor=west] at (-1.2,0.2) {Fase 2: a cada pergunta};
\node[font=\small\itshape] (q) at (-0.4,-1.0) {pergunta};
\node[det,right=5mm of q] (b0) {Primeira\\busca};
\node[llm,right=5mm of b0] (pl) {Planejador};
\node[det,right=5mm of pl] (ex) {Executor};
\node[llm,right=5mm of ex] (re) {Refletor};
\node[det,right=5mm of re] (ct) {Montar\\contexto};
\node[font=\small\itshape,below=5mm of ct] (lt) {leitor};
\draw[arr] (q)--(b0); \draw[arr] (b0)--(pl); \draw[arr] (pl)--(ex); \draw[arr] (ex)--(re); \draw[arr] (re)--(ct); \draw[arr] (ct)--(lt);
\draw[arr,dashed,rounded corners=4pt] (re.south)--++(0,-6mm)-|node[lab,pos=0.25,below]{sem prova ou prova recusada: o motivo volta ao Planejador (até 2 ciclos)}(pl.south);
\draw[arr,dotted] (m.south) -- (ex.north);
\end{tikzpicture}
\caption{As duas fases. Caixas laranja usam o modelo de linguagem (LLM); caixas
azuis são programas comuns, determinísticos. A seta pontilhada indica que a
busca lê a memória construída na fase 1.}
\label{fig:visao}
\end{figure}

\textbf{Fase 1: construir a memória.} Acontece uma vez, antes de qualquer
pergunta. A conversa é guardada de duas formas ligadas entre si: como
\emph{texto} (falas e trechos, com datas), que é o que o leitor vai ler, e como
um \emph{grafo de fatos}, em que as provas são procuradas. Seção~\ref{sec:memoria}.

\textbf{Fase 2: responder.} Para cada pergunta, três papéis trabalham em ciclo:
\begin{itemize}
\item o \textbf{Planejador} (LLM) decide \emph{o que procurar} e escreve um
  plano de busca;
\item o \textbf{Executor} (programa) \emph{segue o plano} pelo grafo e devolve
  as provas que achou, ou o ponto exato em que a prova quebrou;
\item o \textbf{Refletor} (LLM) \emph{confere} a prova lendo as falas
  originais; se não há prova, ou ela não se sustenta, explica o motivo ao
  Planejador, que tenta um novo plano.
\end{itemize}
Ao final, os trechos da prova confirmada entram no contexto do leitor.
Seções~\ref{sec:plano} a~\ref{sec:contexto}.

Essa divisão é a de um investigador: primeiro decide o que precisa descobrir,
depois vai atrás, e antes de concluir confere se o que achou se sustenta. Se
não se sustenta, entende por quê e muda a estratégia.

\begin{porque}
O LLM é bom para \emph{entender} (traduzir uma pergunta em plano, julgar se
uma fala apoia uma resposta) e ruim para \emph{procurar} com garantia em
milhares de fatos. Um programa é o contrário. Por isso o LLM planeja e julga,
e o programa procura. Todo número que decide a busca (similaridades, limiares,
custos) é calculado pelo Executor, e pode ser auditado.
\end{porque}

% =====================================================================
\section{Fase 1: como a memória é construída}
\label{sec:memoria}

A construção transforma a conversa bruta em memória em cinco passos. A
Figura~\ref{fig:grafo} mostra o resultado para a nossa conversa de exemplo.

\subsection{Falas e trechos}

A conversa é uma sequência de \textbf{falas}. Cada fala tem um identificador,
o falante, o texto e a data da sessão em que foi dita. As falas são agrupadas
em \textbf{trechos}: blocos contíguos de até 2048 tokens, sem cortar uma fala
ao meio. Os trechos são a unidade que o leitor lê no final. Um trecho pode
cobrir mais de uma sessão, então guarda o \emph{intervalo} das datas das suas
falas.

\begin{exemplo}[: uma fala]
\texttt{[D4:6] Melanie: I just took my fam camping in the mountains last week!}\\
Sessão de 27/6/2023. Essa fala está no trecho $p_2$, que cobre as sessões de
9/6 a 3/7/2023.
\end{exemplo}

\subsection{Fatos extraídos}

Um LLM lê o texto em janelas curtas e escreve os \textbf{fatos} que encontra,
no formato \emph{(sujeito, relação, objeto)}, mais uma expressão de tempo
quando houver. A extração é adaptada a diálogo: o falante é o sujeito
(``I'' vira ``Melanie''), pronomes são resolvidos dentro da janela, e a
expressão de tempo é copiada como está.

\begin{exemplo}[: da fala ao fato]
Da fala acima sai o fato (Melanie, \rel{camp at}, mountains), com a expressão
de tempo ``last week''.
\end{exemplo}

A relação é \emph{o verbo que o texto usou}: não existe uma lista fixa de
relações. Isso tem uma consequência importante que volta mais tarde: a busca
não pode exigir que a relação do plano seja idêntica à da memória
(``live in'' e ``move to''); ela compara relações por \emph{significado}.

Cada fato guarda também a \textbf{fala de origem}: a fala do trecho que mais
compartilha palavras com o fato (palavras do objeto contam em dobro, as da
relação pela metade). É por essa ligação que, mais tarde, uma prova no grafo
vira texto que o leitor e o verificador podem ler.

\subsection{Datas: quando o evento aconteceu}

A data que importa é a do \emph{evento}, não a da conversa. ``Last week'',
dito em 27/6/2023, significa a semana de 19 a 25/6. A construção resolve a
expressão de tempo usando a data da fala de origem. Quando não há expressão,
o fato recebe o dia da sessão. Toda data é guardada como um
\textbf{intervalo de dias}, porque ``last week'' ou ``in May'' não são um dia só.

\begin{center}\small
\begin{tabular}{@{}llll@{}}
\toprule
Fala dita em & Expressão & Data do evento guardada & Largura \\
\midrule
8/5/2023 & ``yesterday'' & 7/5/2023 & 1 dia \\
27/6/2023 & ``last week'' & 19/6 a 25/6/2023 & 7 dias \\
14/8/2023 & ``last month'' & 1/7 a 31/7/2023 & 31 dias \\
6/7/2023 & (nenhuma) & 6/7/2023, o dia da sessão & 1 dia \\
\bottomrule
\end{tabular}
\end{center}

\subsection{Importância: o quanto o evento foi marcante}

Algumas coisas que alguém conta são corriqueiras (``fui ao mercado''), outras
são marcantes (``foi incrível, estou tão feliz!''). A memória guarda uma nota
de \textbf{importância} entre 0 e 1 para cada fala, e o fato herda a nota da
sua fala de origem. A nota mede a intensidade emocional do texto, contando
palavras de um léxico pequeno e auditável: palavras emotivas comuns valem 1
(``happy'', ``amazing'', ``worried''), intensas valem 1,5 (``thrilled'',
``devastated''), intensificadores valem 0,3 (``so'', ``really'') e cada
exclamação vale 0,5 (até três). A soma é a ``energia'' $E$ da fala, e a nota
satura: a primeira palavra emotiva pesa muito, a décima quase nada.

\begin{formula}
\[
\iota(u)=1-e^{-E(u)/\kappa},\qquad \kappa=2 .
\]
$\iota$ (iota) é a importância da fala $u$. Com $E=0$, $\iota=0$; quanto maior
$E$, mais $\iota$ se aproxima de 1, sem nunca chegar.
\end{formula}

\begin{exemplo}[: notas de importância]
\begin{tabular}{@{}lcc@{}}
``I went to the market.'' & $E=0$ & $\iota=0$ \\
``I'm thrilled!!'' & $E=1{,}5+2\cdot0{,}5=2{,}5$ & $\iota=0{,}71$ \\
``It was amazing, I'm so happy!'' & $E=1+1+0{,}3+0{,}5=2{,}8$ & $\iota=0{,}75$ \\
\end{tabular}
\end{exemplo}

A importância é calculada uma vez e \emph{não depende da pergunta}. Ela só
entra na busca quando o plano pede (``qual foi o momento mais marcante?'').

\subsection{Nomes que são a mesma coisa, relações que são o mesmo verbo}

``Mel'' e ``Melanie'' são a mesma pessoa e devem virar um único nó do grafo.
Parece fácil juntar nomes parecidos, mas é perigoso: para um modelo de
embeddings, ``the first half of the game'' e ``the second half of the game''
são quase idênticos (cosseno 0,95) e não são a mesma coisa.

Por isso dois nomes só são unidos quando as \emph{duas} condições valem:
\begin{itemize}
\item são parecidos em significado (cosseno $\ge0{,}80$); \emph{e}
\item as palavras de um estão contidas nas do outro, com no máximo uma
  palavra a mais para nomes curtos (``Mel'' $\subset$ ``Melanie'',
  ``New York'' $\subset$ ``New York City'').
\end{itemize}
``First'' e ``second'' são palavras que nenhum dos dois lados contém no outro,
então aqueles dois não se unem.

\begin{porque}
Unir dois nomes que são coisas diferentes \emph{fabrica provas falsas}: uma
cadeia poderia passar de uma pessoa para outra. Deixar de unir dois nomes que
são a mesma coisa só \emph{perde} uma prova. O primeiro erro é pior, então a
regra é conservadora.
\end{porque}

Com as relações acontece algo parecido. Formas do mesmo verbo (``paint'',
``painted'', ``painting'') são agrupadas em uma \emph{família} pelo radical de
cada palavra, mas as preposições são mantidas: ``work at'' e ``work for''
continuam diferentes. Relações de famílias diferentes são comparadas por
significado, com embeddings.

\subsection{O grafo}

Com os nomes unidos, cada pessoa ou coisa vira um \textbf{nó}, e cada fato vira
uma \textbf{seta} do sujeito para o objeto, com a relação escrita na seta. A
seta carrega a data do evento, a importância, a confiança da extração e a
ligação com a fala de origem (Figura~\ref{fig:grafo}).

\begin{figure}[H]
\centering
\resizebox{\textwidth}{!}{%
\begin{tikzpicture}[
  ent/.style={draw,ellipse,thick,fill=orange!12,font=\small,inner sep=2pt,minimum height=7mm},
  e/.style={-{Stealth[length=2.2mm]},thick},
  el/.style={font=\scriptsize,inner sep=1pt,align=center},
  pas/.style={draw,rounded corners=2pt,fill=gray!8,font=\scriptsize,align=left,inner sep=4pt,text width=6.2cm}]
\node[ent] (bet) at (-7.2,1.0) {Beth};
\node[ent] (car) at (-3.4,1.0) {Caroline};
\node[ent] (bos) at (-7.2,-1.3) {Boston};
\node[ent] (mel) at (-0.4,0) {Melanie};
\node[ent] (pot) at (-4.0,-1.4) {pottery class};
\node[ent] (mon) at (3.3,1.5) {mountains};
\node[ent] (bea) at (3.7,0) {beach};
\node[ent] (for) at (3.3,-1.5) {forest};
\draw[e] (bet)--node[el,above]{$f_4$: mentor of\\8/5}(car);
\draw[e] (bet)--node[el,left]{$f_5$: move to\\jul/2023}(bos);
\draw[e] (mel)--node[el,above,sloped,pos=0.55]{$f_6$: sign up for, 2/7}(pot);
\draw[e] (mel)--node[el,above,sloped,pos=0.4]{$f_1$: camp at, 19--25/6}(mon);
\draw[e] (mel)--node[el,above]{$f_2$: camp at, 6/7}(bea);
\draw[e] (mel)--node[el,below,sloped]{$f_3$: camp at, 15/7}(for);
\node[pas,anchor=west] at (5.0,0) {%
\textbf{De onde veio cada seta}\\[2pt]
$f_1$: trecho $p_2$, fala D4:6 (27/6): ``I just took my fam camping in the mountains \emph{last week}!''\\[2pt]
$f_2$: trecho $p_3$, fala D6:16 (6/7): ``Here's a pic of my family camping at the beach''\\[2pt]
$f_3$: trecho $p_4$, fala D8:3 (15/7): ``we went camping in the forest''\\[2pt]
$f_4$: trecho $p_1$ (8/5);\ \ $f_5$: trecho $p_5$ (14/8, ``\emph{last month}'');\ \ $f_6$: trecho $p_3$ (2/7)};
\end{tikzpicture}}
\caption{A memória da conversa de exemplo. Cada seta é um fato, com a relação e
a data do evento. À direita, a ligação de cada seta com a fala de onde veio.
(Conversa didática no formato do LoCoMo; as falas são simplificadas.)}
\label{fig:grafo}
\end{figure}

Dois detalhes do grafo importam para a busca:
\begin{itemize}
\item \textbf{As setas têm direção.} ``Beth é mentora de Caroline'' vai de Beth
  para Caroline. A busca respeita a direção: não inverte relações sozinha.
\item \textbf{Pode haver várias setas entre os mesmos nós.} Se Melanie contou
  duas vezes que acampou na praia, em sessões diferentes, são duas setas, cada
  uma com sua data e sua fala. (Em matemática, isso é um \emph{multigrafo}.)
\end{itemize}

A Tabela~\ref{tab:mem} resume o que cada peça da memória guarda.

\begin{table}[H]
\centering\small
\begin{tabularx}{\textwidth}{@{}l Y Y@{}}
\toprule
Peça & O que guarda & Para que serve \\
\midrule
Fala & identificador, falante, texto, data da sessão, importância & ser a prova legível: o verificador e o leitor leem falas \\
Trecho & as falas em ordem, intervalo de datas, importância média, vetor de embedding & ser lido pelo leitor; busca por palavras e por significado \\
Nó & nome principal, variantes unidas (``Mel''), vetor & ancorar os nomes que aparecem no plano \\
Seta (fato) & sujeito, relação, objeto, data do evento, importância, confiança, trecho e fala de origem, vetor da frase & compor a prova \\
\bottomrule
\end{tabularx}
\caption{O que a memória guarda. A ligação seta $\to$ fala é a ponte entre o grafo (onde se prova) e o texto (o que se lê).}
\label{tab:mem}
\end{table}

\begin{formula}[: a memória]
Um \textbf{fato} é a lista de sete campos
\[
f=\big(s_f,\ r_f,\ o_f,\ \tau(f),\ \iota(f),\ c_f,\ \pi(f)\big):
\]
sujeito $s_f$, relação $r_f$, objeto $o_f$, intervalo do evento $\tau(f)$ (tau),
importância $\iota(f)$, confiança da extração $c_f\in[0,1]$ e origem
$\pi(f)=(p_f,u_f)$, o trecho e a fala de onde veio.

O \textbf{grafo} é $\mathcal G=(V,F)$: $V$ são os nós (classes de nomes unidos)
e $F$ são os fatos, cada um uma seta de $[s_f]$ para $[o_f]$, onde $[n]$ é o nó
do nome $n$.

Para comparar relações, $\sem_R(r,r')=1$ se $r$ e $r'$ são da mesma família e,
caso contrário, $\sem_R(r,r')=\cos(\phi(r),\phi(r'))$, onde $\phi$ é o vetor de
embedding (modelo BGE) e $\cos$ o cosseno entre vetores.
\end{formula}

% =====================================================================
\section{O plano: a pergunta escrita na língua do grafo}
\label{sec:plano}

O grafo não entende português nem inglês. Para procurar nele, a pergunta
precisa ser reescrita como um \textbf{padrão de setas}: quais setas procurar,
entre quais nós. Esse padrão é o centro do plano.

\subsection{Átomos e incógnitas}

Cada seta procurada é um \textbf{átomo}, escrito
\rel{relação}(sujeito, objeto). Onde o sujeito ou o objeto é conhecido, vai o
nome (``Melanie''); onde é o que se quer descobrir, vai uma \textbf{incógnita},
marcada com ``?'' ($\V x$, $\V y$). A incógnita $\V x$ é sempre a resposta.

\begin{center}\small
\begin{tabularx}{\textwidth}{@{}>{\raggedright\arraybackslash}p{4.6cm} l Y@{}}
\toprule
Pergunta & Padrão & Forma \\
\midrule
Onde Melanie já acampou? & $\rel{camp at}(\text{Melanie},\V x)$ & um átomo \\
Para que cidade se mudou a mentora de Caroline? & $\rel{mentor of}(\V y,\text{Caroline})\wedge\rel{move to}(\V y,\V x)$ & cadeia: $\V y$ liga os dois átomos \\
Quem é amigo de Ana e trabalha com Bia? & $\rel{friend of}(\V x,\text{Ana})\wedge\rel{work with}(\V x,\text{Bia})$ & interseção: $\V x$ nos dois \\
Quando Melanie acampou em junho? & $\rel{camp at}(\text{Melanie},\V x)\ @\ \V t$ & a resposta é a data $\V t$ do fato \\
\bottomrule
\end{tabularx}
\end{center}

O símbolo $\wedge$ quer dizer ``e'': todos os átomos precisam ser encontrados.
O que dá força ao padrão é a \textbf{incógnita repetida}. Na cadeia, $\V y$
aparece nos dois átomos, e isso obriga a mentora do primeiro fato a ser
\emph{a mesma pessoa} que se mudou no segundo. Sem essa amarra, qualquer
mudança de qualquer pessoa serviria.

\subsection{As outras partes do plano}

Além do padrão, o plano diz que \emph{forma} a resposta tem e \emph{onde no
tempo} procurar:
\begin{itemize}
\item \textbf{Agregação}: um valor só (``qual cidade?''), um conjunto
  (``onde já acampou?'': todos os lugares), uma contagem (``quantas vezes?'').
  Comparações (``mais'', ``menos'') são tratadas como conjunto: o grafo traz os
  fatos a comparar e o leitor compara.
\item \textbf{Tipo da resposta} (opcional): ``martial art'', ``place''. Evita
  que ``o que John pratica?'' case com tudo que John faz.
\item \textbf{Período de referência}: \emph{hoje} (o padrão), \emph{início} da
  conversa (``a primeira vez que\dots'') ou uma \emph{janela} citada (``em
  junho de 2023'').
\item \textbf{Pesos}: quanto o tempo e a importância devem contar na busca,
  em três níveis (nenhum, normal, forte). A próxima seção explica.
\end{itemize}

O plano é escrito pelo Planejador em JSON. Para a pergunta da mentora:

\begin{tcolorbox}[colback=gray!4,colframe=gray!40,boxrule=0.4pt,arc=1pt,left=4pt,right=4pt,top=2pt,bottom=2pt]
\footnotesize\ttfamily
\{"answer\_var": "x",\\
\ "atoms": [\{"relation": "mentor of", "subject": "?y", "object": "Caroline"\},\\
\ \ \ \ \ \ \ \ \ \ \ \{"relation": "move to", \ "subject": "?y", "object": "?x"\}],\\
\ "aggregation": "none",\\
\ "period": \{"reference": "now", "text": ""\},\\
\ "time\_weight": "normal", "importance\_weight": "normal"\}
\end{tcolorbox}

\begin{formula}[: plano e testemunha]
O padrão é uma \emph{consulta conjuntiva} com até 4 átomos:
\[
Q(\V x)\ \leftarrow\ r_1(t_1,t'_1)\ \wedge\ \dots\ \wedge\ r_n(t_n,t'_n)\qquad(n\le4),
\]
em que cada $t_i,t'_i$ é um nome ou uma incógnita. Lê-se: ``$\V x$ é resposta
se existem setas $r_1,\dots,r_n$ ligadas desse jeito''. O plano completo é
$\mathcal P=(Q,\ \alpha,\ \theta,\ I^\star,\ \mathbf w)$: padrão, agregação,
tipo, período e pesos.

Uma \textbf{testemunha} é uma escolha de nós para as incógnitas (chamada
$\beta$) junto com um conjunto de fatos $W=\{f_1,\dots,f_n\}$, um para cada
átomo, tal que cada $f_i$ é uma seta que liga os nós certos com uma relação
parecida com $r_i$, e nenhum pedaço menor de $W$ basta. A resposta é
$\beta(\V x)$, o nó escolhido para $\V x$.
\end{formula}

% =====================================================================
\section{Uma nota para ordenar tudo}
\label{sec:pontuacao}

Em vários momentos a busca precisa ordenar coisas: quais trechos são
melhores, quais setas casam melhor com um átomo. O WitnessRAG usa \emph{uma
única nota} para tudo isso, somando três sinais entre 0 e 1:
\begin{itemize}
\item \textbf{Similaridade}: \emph{isso se parece com o que procuro?} Para um
  trecho, é a busca por palavras (BM25) combinada com a busca por significado
  (embeddings), normalizada entre 0 e 1. Para uma seta, é o quanto ela casa com
  o átomo do plano (explicado na Seção~\ref{sec:executor}).
\item \textbf{Importância}: \emph{foi marcante?} A nota guardada na memória.
\item \textbf{Proximidade no tempo}: \emph{aconteceu perto do período pedido?}
  Vale 1 dentro do período e cai conforme a distância em dias aumenta.
\end{itemize}

\begin{formula}[: a nota]
\[
S(x)=w_{\text{sim}}\cdot\sem(x)+w_{\text{imp}}\cdot\iota(x)+w_{\text{t}}\cdot\prox(x),
\]
\[
\prox(x)=\exp\!\Big(-\frac{\text{dias entre }\tau(x)\text{ e o período}}{\lambda}\Big).
\]
Os pesos $w$ somam 1. A escala $\lambda$ (lambda) diz quão rápido a proximidade
cai: para uma janela citada, é metade da largura da janela (no mínimo 7 dias);
para \emph{hoje} ou \emph{início}, é um quarto da duração da conversa (no
mínimo 7 dias).
\end{formula}

A Figura~\ref{fig:prox} mostra a proximidade para o período ``junho de 2023''.
Dentro de junho vale 1. Seis dias depois do fim de junho vale 0,67; quinze dias
depois, 0,37; um mês depois, 0,14.

\begin{figure}[H]
\centering
\begin{tikzpicture}
\begin{axis}[width=0.8\textwidth,height=4.6cm,xmin=-40,xmax=60,ymin=0,ymax=1.12,
  xlabel={dias em relação a junho de 2023},ylabel={$\prox$},
  xtick={-30,0,30,60},xticklabels={1/5,\ 1/6 a 30/6,30/7,29/8},
  ytick={0,0.5,1},grid=major,grid style={gray!20},axis lines=left,
  label style={font=\small},tick label style={font=\scriptsize}]
\addplot[blue!60!black,very thick,domain=-40:-15,samples=40] {exp((x+15)/15)};
\addplot[blue!60!black,very thick,domain=-15:15] {1};
\addplot[blue!60!black,very thick,domain=15:60,samples=50] {exp(-(x-15)/15)};
\fill[blue!8] (axis cs:-15,0) rectangle (axis cs:15,1);
\node[font=\scriptsize] at (axis cs:0,0.5) {junho};
\addplot[only marks,mark=*,orange!80!black,mark size=2.2pt] coordinates {(8,1) (21,0.670) (30,0.368)};
\node[font=\scriptsize,above] at (axis cs:8,1) {$f_1$};
\node[font=\scriptsize,above right] at (axis cs:21,0.67) {$f_2$ (6/7): 0,67};
\node[font=\scriptsize,above right] at (axis cs:30,0.368) {$f_3$ (15/7): 0,37};
\end{axis}
\end{tikzpicture}
\caption{Proximidade no tempo para o período ``junho de 2023'' ($\lambda=15$
dias). Os pontos são os três acampamentos de Melanie. (O eixo horizontal está
simplificado: junho aparece como a faixa central.)}
\label{fig:prox}
\end{figure}

A mesma fórmula cobre três jeitos de uma pergunta falar de tempo:
\begin{itemize}
\item \textbf{Período = hoje}: favorece o mais recente (``onde Melanie mora
  agora?'').
\item \textbf{Período = início da conversa}: favorece o mais antigo (``qual
  foi o primeiro emprego de Ines?'').
\item \textbf{Período = janela citada}: favorece o que aconteceu naquela data
  (``em junho'').
\end{itemize}

\textbf{Os pesos têm três níveis.} No nível \emph{normal} o tempo e a
importância valem 0 e toda a nota é similaridade. No nível \emph{forte}, o
tempo vale 0,3 e a importância 0,1. O Planejador liga o nível forte quando a
pergunta cita uma data, fala de ``primeira vez'' ou ``mais recente'', ou pede
o que foi mais marcante. A similaridade sempre fica com pelo menos metade do
peso.

\begin{porque}
Os valores vieram de uma calibração sem LLM nas dez conversas do LoCoMo,
medindo se o trecho certo entra entre os cinco escolhidos. Quando a pergunta
não cita data, dar peso ao tempo mudou até 72\% dos contextos \emph{sem}
melhorar nada, por isso o nível normal é zero. Quando a pergunta cita uma data,
peso 0,3 no tempo melhorou em todas as categorias. E a similaridade fica com
pelo menos metade para que um trecho sem relação com a pergunta nunca passe à
frente só por ser recente ou emotivo. Uma consequência útil: com pesos normais,
a ordem é \emph{exatamente} a da busca híbrida comum, então o método nunca
piora o caso fácil por causa da nota.
\end{porque}

% =====================================================================
\section{Fase 2: o ciclo Planejador, Executor, Refletor}
\label{sec:ciclo}

\begin{figure}[H]
\centering
\begin{tikzpicture}[
  box/.style={draw,rounded corners=3pt,thick,minimum height=11mm,minimum width=27mm,font=\sffamily\small,align=center},
  llm/.style={box,fill=orange!15}, det/.style={box,fill=blue!8},
  arr/.style={-{Stealth[length=2.2mm]},thick}, lab/.style={font=\scriptsize,align=center}]
\node[det] (b0) {Primeira busca\\\scriptsize similaridade pura};
\node[llm,right=18mm of b0] (pl) {Planejador\\\scriptsize o que procurar};
\node[det,right=18mm of pl] (ex) {Executor\\\scriptsize segue o plano};
\node[llm,below=22mm of ex] (re) {Refletor\\\scriptsize confere e explica};
\node[det,below=22mm of b0] (ct) {Montar contexto};
\draw[arr] (b0)--node[lab,above]{20 trechos\\(evidências)}(pl);
\draw[arr] (pl)--node[lab,above]{plano}(ex);
\draw[arr] (ex)--node[lab,right]{provas, ou\\onde quebrou}(re);
\draw[arr] (re.south)--++(0,-5mm)-|node[lab,pos=0.25,below]{prova confirmada}(ct.south);
\draw[arr,dashed] (re.west)-|node[lab,pos=0.75,left]{motivo da falha\\+ sonda\\(novo ciclo, até 2)}(pl.south);
\end{tikzpicture}
\caption{O ciclo de busca de uma pergunta.}
\label{fig:ciclo}
\end{figure}

\subsection{Antes do ciclo: a primeira busca}

Antes de planejar, o sistema faz uma busca comum por similaridade e separa os
\textbf{20 trechos} mais parecidos com a pergunta. Eles são as
\textbf{evidências}: não vão direto para o leitor, servem para o Planejador
ver como a memória fala do assunto.

\begin{porque}
Planejar às cegas produz planos com palavras que a memória não usa. Se a
pergunta diz ``live in'' e a memória diz ``move to'', um plano escrito só a
partir da pergunta falha. Mostrando ao Planejador os fatos que a primeira busca
achou, ele pode escrever o plano \emph{com as palavras da memória}.
\end{porque}

\subsection{Planejador: decide o que procurar}

O Planejador é uma chamada ao LLM. Ele recebe:
\begin{itemize}
\item a pergunta;
\item até 30 fatos tirados das evidências, os mais parecidos com a pergunta,
  cada um com a data do evento, no formato
  \texttt{Beth | move to | Boston | July 2023 (chunk 4)};
\item o vocabulário da memória: relações e nomes mais próximos da pergunta;
\item a partir do segundo ciclo, o que deu errado nos ciclos anteriores.
\end{itemize}
E devolve o plano em JSON (Seção~\ref{sec:plano}). Ele \textbf{nunca} vê a
categoria da pergunta no benchmark, a resposta certa ou qualquer anotação. Se o
JSON vier quebrado, o plano é marcado como inválido e o ciclo segue.

\subsection{Executor: segue o plano no grafo}
\label{sec:executor}

O Executor é um programa comum, sem LLM. Ele recebe o plano e trabalha em
quatro passos.

\paragraph{Passo 1: ancorar os nomes.} Cada nome do plano precisa virar um
ponto de partida no grafo. ``Caroline'' casa com o nó Caroline (idêntico,
nota 1,0) e, se houver, com até cinco nós de nome parecido (cosseno
$\ge0{,}55$). Esses nós são as \textbf{âncoras}.

\paragraph{Passo 2: achar candidatos para cada átomo.} Para cada átomo, o
Executor olha as setas que saem ou chegam às âncoras e dá a cada uma uma
\textbf{nota de casamento} $m$, entre 0 e 1. A nota combina três perguntas:
\begin{itemize}
\item \emph{A relação é parecida?} (peso 0,45) Cosseno entre a relação do
  átomo e a da seta, ou 1 se forem da mesma família.
\item \emph{Os nomes batem?} (peso 0,35) A nota da âncora.
\item \emph{A frase inteira é parecida?} (peso 0,20) Cosseno entre
  ``Melanie camp at ?x'' e a frase do fato, ``Melanie camp at beach''.
\end{itemize}
Há dois vetos, que zeram a nota: a relação ser muito diferente (similaridade
abaixo de 0,65), ou o nome estar do lado errado da seta (o plano pede Caroline
como objeto e ela é o sujeito). As 60 setas de nota mais alta são os
\textbf{candidatos} do átomo. Quando o plano liga os pesos de tempo ou
importância, a nota usada para ordenar é a nota única da Seção~\ref{sec:pontuacao},
com $m$ no lugar da similaridade.

\begin{formula}[: nota de casamento]
\[
m(a,f)=\frac{0{,}45\cdot\sem_R(r_a,r_f)\ +\ 0{,}35\cdot s_{\text{nomes}}\ +\ 0{,}20\cdot s_{\text{frase}}}
{0{,}45\ +\ 0{,}35\cdot[\text{o átomo tem nome}]\ +\ 0{,}20}
\]
É uma média ponderada. O denominador só muda quando o átomo não tem nenhum
nome ($\rel{move to}(\V y,\V x)$): aí o peso dos nomes some e a média é refeita
com os outros dois. $m=0$ se $\sem_R<0{,}65$ ou se um nome não casa com o lado
certo da seta.
\end{formula}

\begin{exemplo}[: nota de uma seta]
Átomo $\rel{camp at}(\text{Melanie},\V x)$ e seta $f_2$ = (Melanie, camp at,
beach). Relação idêntica: 1. Nome idêntico: 1. Frases parecidas: 0,82
(valor ilustrativo). Então $m=0{,}45\cdot1+0{,}35\cdot1+0{,}20\cdot0{,}82=0{,}964$.

A seta $f_6$ = (Melanie, sign up for, pottery class) sai do mesmo nó, mas
``sign up for'' não se parece com ``camp at'': abaixo de 0,65, veto, $m=0$.
\end{exemplo}

\paragraph{Passo 3: juntar os átomos (a caminhada).} Agora o Executor monta as
provas. Ele anda \textbf{camada por camada}, um átomo por vez, começando pelo
átomo com mais nomes (o mais restrito, com menos opções). A cada camada, cada
prova parcial tenta crescer com um candidato do próximo átomo, e só cresce se
as incógnitas \textbf{concordam}: se na camada anterior $\V y$ virou Beth, a
seta da camada seguinte tem de sair de Beth. Provas parciais que não concordam
são descartadas. A Figura~\ref{fig:caminhada} mostra a caminhada da pergunta
da mentora.

\begin{figure}[H]
\centering
\begin{tikzpicture}[
  ent/.style={draw,ellipse,thick,fill=orange!12,font=\small,inner sep=2pt,minimum height=7mm},
  e/.style={-{Stealth[length=2.2mm]},thick},
  ok/.style={e,draw=green!50!black,line width=1.4pt},
  no/.style={e,draw=red!70!black,line width=1.2pt,dashed},
  el/.style={font=\scriptsize,align=center},
  stepbox/.style={draw,rounded corners=2pt,fill=white,font=\scriptsize,align=left,text width=5.4cm,inner sep=4pt}]
\node[ent] (car) at (0,0) {Caroline};
\node[ent] (bet) at (-3.3,0) {Beth};
\node[ent] (bos) at (-3.3,-2.2) {Boston};
\draw[ok] (bet)--node[el,above]{$f_4$: mentor of\\$m=0{,}98$}(car);
\draw[no] (bet)--node[el,left,text width=2.3cm,align=right]{$f_5$: move to\\ciclo 1: veto\\ciclo 2: $m=0{,}95$}(bos);
\node[stepbox,anchor=west] at (1.6,0.3) {\textbf{Camada 1} (átomo com nome):\\ procurar setas \rel{mentor of} que \emph{chegam} a Caroline $\Rightarrow f_4$, e $\V y$ vira Beth.};
\node[stepbox,anchor=west] at (1.6,-1.7) {\textbf{Camada 2}: procurar setas que \emph{saem} de Beth ($\V y$ já fixado) com relação parecida com a do plano.\\ Ciclo 1 (``live in''): só há $f_5$, com relação 0,58 $<$ 0,65: \textcolor{red!70!black}{camada vazia, prova quebra aqui}.\\ Ciclo 2 (``move to''): $f_5$ casa: \textcolor{green!40!black}{prova completa}.};
\end{tikzpicture}
\caption{A caminhada para ``Para que cidade se mudou a mentora de Caroline?''. A
incógnita $\V y$ fixada na camada 1 restringe a camada 2 às setas que saem de
Beth.}
\label{fig:caminhada}
\end{figure}

Quando há muitas provas parciais ao mesmo tempo, o Executor guarda as 400
melhores (uma \emph{busca em feixe}). ``Melhor'' é medido por um
\textbf{custo}: a prova é mais barata quando suas setas casam bem e quando
ela usa poucos trechos diferentes.

\begin{formula}[: custo de uma prova]
\[
\cost(W)=-\sum_{f\in W}\log s(f)\ +\ 0{,}15\cdot(\text{número de trechos distintos usados por }W)
\]
$s(f)$ é a nota da seta (entre 0 e 1). O logaritmo transforma notas altas em
custos pequenos: $-\log 0{,}98=0{,}02$, $-\log 0{,}5=0{,}69$. Cada trecho
diferente custa 0,15 a mais, porque cada trecho vai ocupar uma das cinco vagas
do leitor.
\end{formula}

\textbf{Quando a prova quebra.} Se uma camada fica vazia, nenhuma prova
completa existe. O Executor não devolve só ``não achei'': devolve a
\textbf{lacuna}, isto é, qual átomo faltou e o que já estava fixado. No
exemplo: ``faltou \rel{live in} a partir de Beth''. Isso transforma a falha em
informação útil para o próximo ciclo.

\paragraph{Passo 4: a prova é boa o bastante?} Achar uma prova não basta; ela
precisa ser \emph{confiável} e \emph{específica}. Cada prova recebe um
\textbf{suporte}: o produto das notas das suas setas, vezes a confiança da
extração de cada fato. Por ser um produto, a prova vale o seu elo mais fraco:
uma seta ruim derruba a prova inteira.

\begin{formula}[: suporte]
\[
\supp(W)=\prod_{f\in W} m(a_f,f)\cdot c_f
\]
$a_f$ é o átomo que o fato $f$ realiza e $c_f$ a confiança da extração (0,9 por
padrão).
\end{formula}

As regras de aceitação dependem da forma da resposta:
\begin{itemize}
\item \textbf{Resposta de um valor} (``qual cidade?''): no máximo 3 respostas
  diferentes e 3 provas, nenhum corte na busca, a melhor prova com suporte
  $\ge0{,}45$, e os trechos novos que ela traz precisam caber no contexto
  (Seção~\ref{sec:contexto}).
\item \textbf{Resposta em conjunto} (``onde já acampou?''): muitas respostas
  são esperadas, então o teste é item a item: cada item precisa da própria
  prova com suporte $\ge0{,}45$. Se o plano tem tipo, os 12 melhores itens
  seguem para o Refletor; sem tipo e com mais de 12 itens, o plano é
  considerado vago demais.
\end{itemize}

\begin{porque}
Se um plano de ``um valor'' encontra 8 respostas diferentes, o problema não é
falta de prova: o plano está vago (``does(John, ?x)'' casa com tudo o que John
faz). Aceitar isso encheria o contexto de ruído. Recusar e explicar o motivo
ao Planejador (``8 respostas, seja mais específico'') leva a um plano melhor.
\end{porque}

\subsection{Refletor: confere a prova e explica a falha}

O Refletor tem dois papéis.

\paragraph{Verificar.} Quando o Executor acha uma prova que vai trazer um
trecho novo para o contexto, o LLM lê a pergunta, o plano e, para cada
resposta candidata, os fatos da prova \emph{junto com as falas originais}
(com falante, data e as falas vizinhas). Por exemplo:

\begin{tcolorbox}[colback=gray!4,colframe=gray!40,boxrule=0.4pt,arc=1pt,left=4pt,right=4pt,top=2pt,bottom=2pt]
\footnotesize\ttfamily
A1: Boston\\
\ \ fact: Beth | mentor of | Caroline (event date: 8 May 2023)\\
\ \ \ \ [D1:5] Caroline: My mentor Beth has helped me so much...\\
\ \ fact: Beth | move to | Boston (event date: July 2023)\\
\ \ \ \ (session date: 14 August 2023)\\
\ \ \ \ [D12:8] Caroline: Beth moved to Boston last month...
\end{tcolorbox}

Para cada resposta, ele decide: \emph{apoiada}, ou \emph{recusada} por um de
quatro motivos (entidade errada, período errado, tipo errado, falta de apoio).
O verificador \emph{julga}, não procura; na dúvida, aceita. Se a prova já está
inteira nos trechos que o leitor ia receber de qualquer jeito, não há o que
verificar e a chamada é dispensada.

\begin{porque}
A extração de fatos erra, e a similaridade entre relações é aproximada. O
grafo pode ``provar'' algo que o texto não diz. Ler a fala original pega esses
erros antes que eles mudem o contexto. Por que mostrar a fala e não só o fato?
Porque o fato é um resumo; a fala tem o contexto (quem disse, quando, sobre o
quê).
\end{porque}

\paragraph{Diagnosticar.} Se não houve prova, ou ela foi recusada, o Refletor
transforma o motivo em uma instrução para o próximo plano:

\begin{center}\small
\begin{tabularx}{\textwidth}{@{}l Y@{}}
\toprule
O que aconteceu & O que o Planejador ouve no próximo ciclo \\
\midrule
Uma camada ficou vazia (lacuna) & nenhum fato casou com \rel{live in} para Beth \\
Respostas demais & o plano casou com 8 respostas diferentes; seja mais específico \\
Busca cortada & o plano é amplo demais; seja mais específico \\
Só casamentos fracos & use as palavras de relação que a memória usa \\
Prova não cabe no contexto & a prova precisa de trechos demais; use menos fatos \\
Verificador recusou & a prova foi recusada (período errado, entidade errada\dots) \\
\bottomrule
\end{tabularx}
\end{center}

Quando há lacuna, ela vira também uma \textbf{sonda}: um texto curto (``Beth
live in'') usado para escolher quais fatos mostrar ao Planejador no próximo
ciclo. Assim o Planejador passa a ver os fatos sobre Beth, entre eles
``Beth | move to | Boston''.

\subsection{Quando o ciclo termina}

O ciclo roda no máximo \textbf{2 vezes} e termina antes se:
\begin{itemize}
\item uma prova foi confirmada (sucesso);
\item o Planejador repetiu um plano já tentado (não há o que ganhar).
\end{itemize}
Se terminar sem prova, o leitor recebe os cinco melhores trechos da busca comum:
o sistema cai para o comportamento da busca híbrida, sem prejuízo. No pior caso,
uma pergunta custa 4 chamadas ao LLM além do leitor (2 planos e 2 verificações).

\begin{formula}[: o ciclo em pseudocódigo]
\begin{tabbing}
\quad\=\quad\=\quad\=\kill
evidências $\leftarrow$ 20 melhores trechos por similaridade;\ \ feedback $\leftarrow$ vazio\\
\textbf{para} ciclo $=1,2$:\\
\> plano $\leftarrow$ \op{Planejador}(pergunta, fatos das evidências, vocabulário, feedback)\\
\> \textbf{se} plano inválido: feedback += erro; \textbf{continue}\\
\> \textbf{se} plano repetido: \textbf{pare}\\
\> resultado $\leftarrow$ \op{Executor}(plano)\\
\> \textbf{se} resultado tem prova aceitável:\\
\> \> prova $\leftarrow$ \op{Refletor}.verificar(pergunta, plano, resultado)\\
\> \> \textbf{se} prova confirmada: \textbf{devolva} contexto com a prova\\
\> feedback += \op{Refletor}.diagnosticar(resultado)\quad(motivo e sonda)\\
\textbf{devolva} os 5 melhores trechos da busca comum
\end{tabbing}
\end{formula}

% =====================================================================
\section{Montando o contexto do leitor}
\label{sec:contexto}

O leitor recebe $k=5$ trechos. O ponto de partida são os cinco melhores pela
nota do plano (com pesos normais, os mesmos da busca híbrida). Uma prova
confirmada muda esse contexto de um jeito controlado
(Figura~\ref{fig:contexto}):
\begin{itemize}
\item os trechos da prova que já estão entre os cinco ficam onde estão;
\item os que faltam entram \textbf{pelo fim da lista}, no lugar dos trechos
  menos bem colocados;
\item no máximo $k_W$ trocas: 2 para planos compostos (cadeia, interseção ou
  conjunto) e 1 para um fato com um valor;
\item as primeiras $\max(1,k-k_W)$ posições nunca mudam.
\end{itemize}

\begin{figure}[H]
\centering
\begin{tikzpicture}[slot/.style={draw,thick,minimum width=17mm,minimum height=8mm,font=\small},
  prot/.style={slot,fill=blue!10},free/.style={slot,fill=white},newp/.style={slot,fill=green!18}]
\node[font=\small,anchor=east] at (-0.3,0) {antes};
\foreach \i/\t/\s in {1/$p_1$/prot,2/$p_8$/prot,3/$p_3$/prot,4/$p_9$/free,5/$p_{12}$/free}
  \node[\s] at (\i*1.9,0) {\t};
\node[font=\small,anchor=east] at (-0.3,-1.3) {depois};
\foreach \i/\t/\s in {1/$p_1$/prot,2/$p_8$/prot,3/$p_3$/prot,4/$p_9$/free,5/$p_5$/newp}
  \node[\s] at (\i*1.9,-1.3) {\t};
\draw[decorate,decoration={brace,amplitude=5pt,mirror}] (1.1,-1.85)--(6.5,-1.85) node[midway,below=5pt,font=\scriptsize]{protegidas: $\max(1,5-2)=3$};
\draw[decorate,decoration={brace,amplitude=5pt,mirror}] (6.9,-1.85)--(10.3,-1.85) node[midway,below=5pt,font=\scriptsize]{podem ser trocadas ($k_W=2$)};
\draw[-{Stealth},thick,green!40!black] (9.5,-0.45)--(9.5,-0.85);
\end{tikzpicture}
\caption{A prova da mentora usa os trechos $p_1$ (já no contexto) e $p_5$
(fora). $p_5$ entra no lugar do último trecho.}
\label{fig:contexto}
\end{figure}

\begin{porque}
A prova pode estar errada mesmo depois de verificada. Proteger as primeiras
posições garante que o método \emph{só acrescenta} evidência provada: nunca
tira do leitor o que a busca comum achou de melhor. Se não houver prova
confirmada, o contexto é exatamente o da busca híbrida.
\end{porque}

\textbf{Opções da versão 4} (desligadas por padrão, usadas nos testes mais
recentes):
\begin{itemize}
\item \textbf{Entrega como falas}: em vez de trocar trechos, as provas que não
  cabem entram como um bloco curto (até 2400 caracteres) com as falas de
  origem, depois dos cinco trechos. Só diálogo, sem os fatos extraídos nem a
  resposta, para o leitor não copiar a extração.
\item \textbf{Premissas para hipóteses}: para perguntas como ``Would Caroline
  likely have Dr.\ Seuss books?'', não existe prova: a resposta depende de
  fatos \emph{e} de conhecimento de mundo. O sistema entrega até 8 falas sobre
  Caroline, escolhidas pelos conceitos que apoiariam ou contradiriam a
  hipótese (``children's books'', ``reading habits''). São premissas para o
  leitor raciocinar, não uma prova.
\end{itemize}

% =====================================================================
\section{Três perguntas do começo ao fim}
\label{sec:exemplos}

Os três exemplos usam a memória da Figura~\ref{fig:grafo}. Os limiares e pesos
são os do código; as similaridades (cossenos) são valores ilustrativos; a
confiança de todos os fatos é 0,9.

\needspace{10\baselineskip}\subsection*{Exemplo A: ``Onde Melanie já acampou?'' (resposta em conjunto)}

\begin{enumerate}
\item \textbf{Primeira busca.} Os trechos mais parecidos com a pergunta
  incluem $p_2$ e $p_3$ (montanhas e praia). O trecho da floresta, $p_4$,
  ficou em 7.º lugar: fora dos cinco.
\item \textbf{Planejador.} Vê entre as evidências fatos como
  ``Melanie | camp at | mountains''. Escreve:
  $\rel{camp at}(\text{Melanie},\V x)$, agregação \emph{conjunto}, tipo
  ``place'', pesos normais.
\item \textbf{Executor.}
  \begin{itemize}
  \item Âncora: o nó Melanie.
  \item Candidatos: as quatro setas que saem de Melanie. $f_1$, $f_2$ e $f_3$
    casam com $m\approx0{,}96$; $f_6$ (\rel{sign up for}) é vetada.
  \item Caminhada: um átomo só, então cada candidato já é uma prova completa.
    Três provas: \{$f_1$\} $\to$ mountains, \{$f_2$\} $\to$ beach,
    \{$f_3$\} $\to$ forest.
  \item Teste: conjunto, então item a item. Cada suporte é
    $0{,}96\times0{,}9\approx0{,}87\ge0{,}45$. Aceito.
  \end{itemize}
\item \textbf{Refletor.} A prova da floresta traz um trecho novo ($p_4$),
  então há verificação. O LLM lê as três falas e confirma os três itens.
\item \textbf{Contexto.} $p_4$ entra no lugar do 5.º trecho. O leitor recebe as
  três viagens e responde ``mountains, beach and forest''.
\end{enumerate}
\emph{O que este exemplo mostra}: a resposta espalhada por várias sessões é
reunida porque o plano pede \emph{todas} as setas \rel{camp at} de Melanie, e
cada uma aponta para o seu trecho.

\needspace{10\baselineskip}\subsection*{Exemplo B: ``Quando Melanie acampou em junho?'' (o tempo decide)}

\begin{enumerate}
\item \textbf{Planejador.} A pergunta cita um mês, então o plano usa
  período \emph{janela} ``June 2023'' e tempo \emph{forte}:
  $\rel{camp at}(\text{Melanie},\V x)\ @\ \V t$, resposta $\V t$ (a data),
  pesos $(0{,}7;\ 0;\ 0{,}3)$.
\item \textbf{Executor.} As três setas \rel{camp at} casam igualmente bem
  ($m\approx0{,}96$). Sem o tempo, elas empatariam. Com o tempo:
\end{enumerate}
\begin{center}\small
\begin{tabular}{@{}lccccc@{}}
\toprule
Seta & data do evento & dias fora de junho & $\prox$ & nota $=0{,}7\cdot m+0{,}3\cdot\prox$ \\
\midrule
$f_1$ & 19--25/6 & 0 & 1,000 & \textbf{0,972} \\
$f_2$ & 6/7 & 6 & 0,670 & 0,873 \\
$f_3$ & 15/7 & 15 & 0,368 & 0,782 \\
\bottomrule
\end{tabular}
\end{center}
\begin{enumerate}\setcounter{enumi}{2}
\item \textbf{Teste.} Resposta de um valor: vale a prova mais bem colocada,
  $f_1$. Três respostas no máximo, suporte alto: aceito.
\item \textbf{Refletor e contexto.} O trecho $p_2$ já estava entre os cinco
  (os mesmos pesos de tempo também reordenam os trechos, e $p_2$ sobe), então
  não há verificação nem troca.
\item \textbf{Leitor.} Responde ``the week before 27 June 2023''.
\end{enumerate}
\emph{O que este exemplo mostra}: a data certa existe porque a construção da
memória resolveu o ``last week'' da fala. E é a proximidade no tempo, não a
similaridade, que separa junho de julho.

\needspace{10\baselineskip}\subsection*{Exemplo C: ``Para que cidade se mudou a mentora de Caroline?'' (cadeia com reflexão)}

\needspace{8\baselineskip}\textbf{Ciclo 1.}
\begin{enumerate}
\item \textbf{Primeira busca.} Traz trechos sobre Caroline, entre eles $p_1$
  (Beth é mentora). O trecho $p_5$, em que Beth se muda, não cita Caroline e
  fica fora dos cinco (mas está entre os 20 das evidências).
\item \textbf{Planejador.} Escreve a cadeia com o verbo da pergunta:
  $\rel{mentor of}(\V y,\text{Caroline})\wedge\rel{live in}(\V y,\V x)$.
\item \textbf{Executor.}
  \begin{itemize}
  \item Camada 1 (o átomo com nome): setas \rel{mentor of} que chegam a
    Caroline. $f_4$ casa com $m=0{,}98$. $\V y$ vira Beth.
  \item Camada 2: setas que saem de Beth com relação parecida com
    \rel{live in}. A única é $f_5$ (\rel{move to}), mas a similaridade entre
    ``live in'' e ``move to'' é 0,58, abaixo de 0,65: veto. Camada vazia.
  \item Lacuna: faltou \rel{live in} a partir de Beth.
  \end{itemize}
\item \textbf{Refletor (diagnóstico).} Escreve ao Planejador ``nenhum fato
  casou com \rel{live in} para Beth'' e cria a sonda ``Beth live in''. Com a
  sonda, os fatos mostrados no próximo ciclo incluem
  ``Beth | move to | Boston | July 2023''.
\end{enumerate}

\textbf{Ciclo 2.}
\begin{enumerate}\setcounter{enumi}{4}
\item \textbf{Planejador.} Vê como a memória fala da mudança e reescreve:
  $\rel{mentor of}(\V y,\text{Caroline})\wedge\rel{move to}(\V y,\V x)$
  (o JSON da Seção~\ref{sec:plano}).
\item \textbf{Executor.} Camada 1 igual. Camada 2: agora $f_5$ casa. O átomo
  \rel{move to}($\V y$, $\V x$) não tem nome, então a média é refeita sem o peso
  dos nomes:
  $m(f_5)=(0{,}45\cdot1+0{,}20\cdot0{,}85)/(0{,}45+0{,}20)=0{,}954$.
  Prova: $W=\{f_4,f_5\}$, $\V y=$ Beth, $\V x=$ Boston.
  \begin{itemize}
  \item suporte $=0{,}98\times0{,}954\times0{,}9\times0{,}9=0{,}76\ge0{,}45$;
  \item custo $=-\log0{,}98-\log0{,}954+0{,}15\times2=0{,}02+0{,}05+0{,}30=0{,}37$.
  \end{itemize}
  Uma resposta, uma prova: aceito.
\item \textbf{Refletor (verificação).} $p_5$ é novo, então o LLM lê as duas
  falas (a caixa da Seção~\ref{sec:ciclo}) e confirma.
\item \textbf{Contexto.} $p_5$ entra na 5.ª posição (Figura~\ref{fig:contexto}).
  O leitor responde ``Boston''.
\end{enumerate}
\emph{O que este exemplo mostra}: (i) a cadeia encontra um trecho que a
similaridade nunca traria, porque ``Boston'' não aparece perto de
``Caroline''; (ii) a falha do primeiro ciclo não foi um beco sem saída: a
lacuna disse exatamente \emph{onde} a prova quebrou, e o segundo plano
corrigiu só isso.

% =====================================================================
\section{O que o método garante e o que não garante}
\label{sec:garantias}

\textbf{Garante:}
\begin{itemize}
\item \textbf{Nunca piora o contexto sem prova.} Sem prova confirmada e com
  pesos normais, o leitor recebe exatamente os trechos da busca híbrida. Com
  prova, no máximo $k_W$ trechos mudam, e as primeiras posições ficam.
\item \textbf{Não usa informação do benchmark.} Nenhuma etapa lê a categoria da
  pergunta, a resposta ou anotações.
\item \textbf{Custo limitado.} No máximo 4 chamadas ao LLM por pergunta além do
  leitor; o Executor é só índice e contas.
\item \textbf{Auditável.} Toda prova aceita tem os fatos, as notas e as falas de
  origem registrados.
\item \textbf{Correção no modo exato.} Se as relações e os nomes forem
  comparados por igualdade (não por similaridade) e nenhum corte ocorrer, o
  Executor encontra exatamente todas as respostas que o grafo prova.
\end{itemize}

\textbf{Não garante:}
\begin{itemize}
\item \textbf{Que os fatos estejam certos.} A prova é tão boa quanto a
  extração. A confiança $c_f$ é um valor fixo, não calibrado.
\item \textbf{Que nada ficou de fora.} Um conjunto contém o que a memória
  prova; um lugar que a extração não registrou não aparece.
\item \textbf{Exatidão no modo padrão.} Comparar relações por similaridade é uma
  aproximação: ``live in'' e ``move to'' podem ou não casar, dependendo do
  limiar.
\item \textbf{Sinônimos sem palavras em comum.} ``TV'' e ``television'' não são
  unidos, porque a regra de nomes é conservadora.
\end{itemize}

% =====================================================================
\section{Glossário}
\label{sec:glossario}

\begin{center}\small
\begin{tabularx}{\textwidth}{@{}l Y@{}}
\toprule
Termo & Significado \\
\midrule
Fala & uma intervenção de um falante na conversa, com a data da sessão \\
Trecho & bloco contíguo de falas (até 2048 tokens); é o que o leitor lê \\
Fato & (sujeito, relação, objeto) extraído de uma fala, com data do evento, importância, confiança e origem \\
Nó & uma pessoa ou coisa do grafo; nomes unidos (``Mel'', ``Melanie'') são um só nó \\
Seta & um fato visto no grafo: vai do sujeito para o objeto \\
Leitor & o LLM que escreve a resposta final a partir de 5 trechos \\
Evidências & os 20 trechos da primeira busca, mostrados ao Planejador \\
Plano & a pergunta reescrita como padrão de setas, com agregação, tipo, período e pesos \\
Átomo & uma seta procurada: \rel{relação}(sujeito, objeto) \\
Incógnita & o que se quer descobrir num átomo ($\V x$, $\V y$); $\V x$ é a resposta \\
Âncora & nó do grafo que corresponde a um nome do plano \\
Nota de casamento ($m$) & quanto uma seta corresponde a um átomo (relação, nomes, frase) \\
Testemunha / prova & o menor conjunto de fatos do qual a resposta decorre \\
Suporte & produto das notas e confianças dos fatos de uma prova \\
Custo & o que a caminhada minimiza: notas baixas e trechos demais custam mais \\
Lacuna & o átomo em que a prova quebrou, com o que já estava fixado \\
Sonda & texto curto feito da lacuna, usado para mostrar novos fatos ao Planejador \\
Período de referência & hoje, início da conversa ou uma janela citada \\
Proximidade & 1 dentro do período, caindo com a distância em dias \\
$k$, $k_W$ & trechos entregues ao leitor (5); trechos que uma prova pode trocar (1 ou 2) \\
\bottomrule
\end{tabularx}
\end{center}

% =====================================================================
\section{Onde cada peça está no código}
\label{sec:codigo}

\begin{center}\small
\begin{tabularx}{\textwidth}{@{}l Y@{}}
\toprule
Peça & Arquivo e função \\
\midrule
Extração de fatos & \texttt{wrag/ie.py} (\texttt{extract\_corpus}) \\
Datas, importância, fala de origem & \texttt{wrag/witness/dated\_memory.py} (\texttt{DatedMemory}, \texttt{arousal}) \\
Grafo, nomes unidos, famílias de relações & \texttt{wrag/graph.py} (\texttt{build\_graph}, \texttt{relation\_signature}) \\
Nota única e proximidade & \texttt{wrag/witness/scoring.py} (\texttt{MemoryScorer}, \texttt{proximity}) \\
Planejador & \texttt{wrag/witness/plan.py} (\texttt{plan\_question}); prompt em \texttt{wrag/prompts.py} \\
Executor: âncoras, candidatos, caminhada & \texttt{wrag/witness/search.py} (\texttt{match\_entity}, \texttt{ground}, \texttt{join}) \\
Executor: suporte e teste de aceitação & \texttt{wrag/methods/witnessrag.py} (\texttt{\_prove}, \texttt{\_prove\_items}) \\
Refletor: verificação & \texttt{wrag/witness/confirm.py} (\texttt{confirm}) \\
Refletor: diagnóstico & \texttt{wrag/methods/witnessrag.py} (\texttt{\_OUTCOME\_FEEDBACK}) \\
Ciclo inteiro e contexto & \texttt{wrag/methods/witnessrag.py} (\texttt{\_retrieve\_proof}) \\
Parâmetros & \texttt{wrag/config.py} (\texttt{WitnessConfig}) \\
\bottomrule
\end{tabularx}
\end{center}

Observação: no código, o Refletor não é uma classe única; seus dois papéis
estão em lugares diferentes (tabela acima). A divisão Planejador, Executor e
Refletor deste guia é a forma de explicar o que o código faz.

\subsection*{Parâmetros principais}

\begin{center}\small
\begin{tabular}{@{}lll@{}}
\toprule
Parâmetro & Valor & Onde aparece \\
\midrule
Trechos para o leitor ($k$) & 5 & contexto \\
Evidências & 20 trechos, 30 fatos & Planejador \\
Ciclos & 2 & ciclo \\
Átomos por plano & até 4 & plano \\
Pesos da nota de casamento & 0,45 / 0,35 / 0,20 & Executor, passo 2 \\
Veto de relação & 0,65 & Executor, passo 2 \\
Nomes: cosseno mínimo / âncoras & 0,55 / até 5 & Executor, passo 1 \\
Candidatos por átomo & 60 & Executor, passo 2 \\
Feixe & 400 & Executor, passo 3 \\
Custo por trecho & 0,15 & Executor, passo 3 \\
Suporte mínimo & 0,45 & Executor, passo 4 \\
Máx. respostas / provas (um valor) & 3 / 3 & Executor, passo 4 \\
Máx. itens de conjunto & 12 & Executor, passo 4 \\
Peso forte: tempo / importância & 0,3 / 0,1 & nota única \\
Trocas no contexto ($k_W$) & 2 composto / 1 simples & contexto \\
União de nomes & cosseno $\ge0{,}80$ + palavras contidas & memória \\
Importância ($\kappa$) & 2 & memória \\
\bottomrule
\end{tabular}
\end{center}

\end{document}
